% ============================================================
%  07_reconstruction.tex — Section 7 : Reconstruction 3D + validation
% ============================================================
\section{Reconstruction 3D et validation (résumé)}

\subsection{Pipeline de reconstruction}

Chaque solution CSP (assignation de tailles) est reconstruite en une molécule 3D selon
trois phases :

\begin{enumerate}
  \item \textbf{Phase A — Modifications topologiques} : pour chaque hexagone dont la taille
    assignée diffère de 6, le cycle est modifié (un sommet retiré pour un pentagone, un sommet
    ajouté pour un heptagone). Seuls les côtés libres ($\sigma_i = 0$) sont touchés.

  \item \textbf{Phase B — Placement géométrique} : un parcours BFS sur le graphe dual place
    les cycles comme des polygones réguliers fusionnés. L'hexagone de plus haut degré est
    choisi comme racine, puis chaque cycle enfant est placé par rapport au parent via l'arête
    partagée.

  \item \textbf{Phase C — Assemblage moléculaire} : les coordonnées sont injectées dans un
    \texttt{MolecularGraph}, les doubles liaisons sont placées par couplage maximum (Blossom),
    les hydrogènes sont ajoutés, et le tout est exporté en XYZ.
\end{enumerate}

Le détail de chaque phase est documenté dans \textit{validation\_reconstruction\_csp.tex}.

\subsection{Validation xTB + planarité}

Chaque XYZ reconstruit est :
\begin{enumerate}
  \item Optimisé par xTB (GFN2-xTB, \texttt{--opt tight})
  \item Testé en planarité par ACP (angle maximal $\leq 10°$)
\end{enumerate}

\subsubsection{Perturbation des coordonnées initiales}

La reconstruction place tous les atomes dans le plan $z = 0$ (géométrie de polygones
réguliers). Or, xTB utilise un algorithme de descente de gradient pour minimiser l'énergie :
si la géométrie initiale est parfaitement plane, les forces perpendiculaires au plan sont
nulles par symétrie, et xTB \textbf{reste piégé dans le minimum local plat}.

Ce problème a été mis en évidence en comparant deux approches pour tester la planarité d'un
même benzénoïde à 5 hexagones :
\begin{itemize}
  \item Placement par coordonnées du réseau hexagonal $\to$ xTB $\to$ \textbf{non plan} (58.6°)
  \item Placement par BFS polygones réguliers $\to$ xTB $\to$ \textbf{plan} (0.0°)
\end{itemize}

Les deux géométries initiales étaient planes ($z = 0$), mais les légères différences de
coordonnées faisaient que xTB convergeait vers des minima différents. La solution adoptée
est d'ajouter une \textbf{perturbation aléatoire} de $\pm 0.1$~\AA{} sur les coordonnées $z$
de chaque atome avant l'optimisation :

\begin{lstlisting}
def _write_perturbed_xyz(input_path, output_path, amplitude=0.1):
    atoms = _read_atoms_from_xyz(input_path)
    with open(output_path, 'w') as f:
        f.write(f"{len(atoms)}\n")
        f.write("perturbed for xTB optimization\n")
        for elem, x, y, z in atoms:
            z_perturbed = z + random.uniform(-amplitude, amplitude)
            f.write(f"{elem} {x:.6f} {y:.6f} {z_perturbed:.6f}\n")
\end{lstlisting}

Si la molécule est réellement plane, xTB la ramène à $z \approx 0$ malgré la perturbation.
Si elle ne l'est pas, xTB trouve le vrai minimum non plan. C'est l'équivalent numérique de la
procédure manuelle suggérée par le chimiste : \textit{``déplacer un atome pendant l'optimisation
pour vérifier s'il revient à sa place''}.

\subsubsection{Inclusion du benzénoïde comme solution}

Le benzénoïde d'origine (toutes les tailles à 6) est désormais inclus dans les solutions du
CSP. Il passe par le même pipeline de reconstruction et de validation xTB que les solutions
non hexagonales. Cela permet de vérifier si le benzénoïde lui-même est planaire — ce qui n'est
pas garanti pour les grandes structures.

\subsubsection{Résultats}

\textbf{Résultats :}
\begin{itemize}
  \item \texttt{first.graph} : \textbf{2/2 solutions planes} (dont le benzénoïde)
  \item \texttt{second.graph} : \textbf{0/41 solutions planes}
\end{itemize}

\begin{remarque}
  Le fait que les 41 solutions de \texttt{second.graph} soient toutes non planes suggère que
  les contraintes locales (table de voisinage) ne suffisent pas pour les grands benzénoïdes.
  Un déficit angulaire cumulatif se crée lorsque trop de pentagones et heptagones sont présents
  dans une structure étendue. La question de contraintes globales (courbure de Gauss discrète,
  contraintes de planéité) reste ouverte et sera discutée avec le chimiste et le tuteur.
\end{remarque}
